\section{Introduction}
Modern sequence-to-sequence foundation models, such as OpenAI Whisper \cite{radford2023robust} and PhoWhisper \cite{phan2022phowhisper}, have fundamentally transformed automatic speech recognition across diverse global languages. However, their transcription accuracy suffers a dramatic drop when exposed to singing voice---even when sung lyrics are linguistically identical to standard conversational utterances \cite{gupta2022deep,zhuo2023lyricwhiz}. This persistent performance gap between speech and song is formally designated as the \textit{Speech-to-Song Mismatch} \cite{ou2022transfer,mesaros2010automatic}.

In tonal languages such as Vietnamese, Mandarin, and Cantonese, this mismatch poses severe linguistic and computational challenges \cite{pham2003vietnamese,nguyen1997vietnamese,wong2002how}. Unlike non-tonal languages where fundamental frequency primarily conveys pragmatic affect or sentential intonation, Vietnamese assigns core lexical contrasts to six canonical tones: \textit{Ngang} (Level), \textit{Huy\`{e}n} (Falling), \textit{S\'{a}c} (Rising), \textit{H\H{o}i} (Dipping), \textit{Ng\~{a}} (Glottalized Broken), and \textit{N\d{a}ng} (Constricted Drop) \cite{brunelle2009tone,kirby2016tone}. These tones are phonologically distinguished by dynamic fundamental frequency contours ($F_0(t)$), duration, and phonation registers (creaky voice, glottal closure). 

When an individual sings, the musical melody dictates vocal cord vibration according to discrete musical scale degrees (e.g., C-E-G), overriding natural lexical pitch contours \cite{kirby2016tone}. This acoustic collision produces \textit{Melodic Tone Neutralization}, wherein syllables with rising or falling lexical tones are flattened into level melodic sustains, causing standard ASR acoustic decoders to commit severe phoneme and tone misclassifications \cite{vu2024vietsing,hoang2026vietprism}.

Concurrently, singing voice features prolonged vowel durations (stretching up to $300\text{--}500\%$ relative to conversational speech) and slow periodic vibrato pitch modulations ($5\text{--}8$\,Hz) \cite{mesaros2010automatic,berendes2024lyrics}. Under standard Transformer cross-attention architectures \cite{vaswani2017attention,bahdanau2015neural}, these sustained acoustic frames cause cross-attention matrices to stall horizontally, as the autoregressive decoder repeatedly queries identical acoustic states \cite{stoller2024jamalt}. Unable to detect expected phonetic transitions, the decoder degenerates into runaway repetition loops, inflating Word Error Rates (WER) past several hundred percent due to spurious token insertions \cite{morris2004wer,levenshtein1966binary}.

To overcome these structural barriers without the prohibitive cost of multi-thousand-hour pre-training from scratch, we propose an integrated multi-phase adaptation and streaming framework tailored for tonal singing voice:
\begin{enumerate}
    \item \textbf{Glottal Inverse Filtering (GIF):} Using Iterative Adaptive Inverse Filtering (IAIF) \cite{alku1992glottal,iseli2004f0}, we decompose the vocal tract filter and extract four glottal phonation parameters (Open Quotient $O_q$, formant-corrected spectral tilt $H_1^*-H_2^*$, Normalized Amplitude Quotient NAQ \cite{alku2002normalized}, and period-doubling subharmonic ratio). We prove empirically that singers preserve creaky glottalization cues ($87.6\%$ retention) for \textit{Ng\~{a}} and \textit{N\d{a}ng} tones even when melodic pitch $F_0$ is completely flattened.
    \item \textbf{Continuous Pitch Conditioning and InfoNCE:} We introduce a continuous pitch embedding block mapping F0 into semitone cents ($256$ bins) fused with learnable acoustic gating, optimized via Syllable-Tone InfoNCE contrastive loss \cite{oord2018infonce,zhou2023rmvpe} to enforce musical pitch invariance.
    \item \textbf{Multi-Dialect Mixture-of-Experts (MoE):} Recognizing deep regional phonological divergences in Vietnamese vocal music, we implement a sparsely-gated LoRA Mixture-of-Experts architecture \cite{hu2021lora,shazeer2017outrageously,fedus2022switch} with an Acoustic Dialect Router dynamic gating network across Northern (Hanoi), Central (Hue/Da Nang), and Southern (Saigon) regional experts.
    \item \textbf{Macro-Scale Pretraining on VietLyrics:} We build a distributed data parallel (DDP) pretraining pipeline utilizing the 647-hour VietLyrics dataset \cite{vietlyrics2025} and our champion multi-component objective ($RS_{01}$), driving the multi-domain Pareto score to a record $26.41\%$.
    \item \textbf{Streaming MoChA Engine:} We design a Monotonic Chunkwise Attention (MoChA) \cite{chiu2018monotonic} streaming decoder operating on $160$\,ms chunks. Tested on 58 real audio files on commodity laptop hardware, our streaming engine runs at $8.50$\,ms chunk latency ($RTF = 0.1091$, $9.2\times$ faster than real-time) while preserving $98.72\%$ of offline recognition accuracy ($55.40\%$ singing WER).
\end{enumerate}
